\section{Related Work}
\label{sec:related_work}

\paragraph{Distance Underestimation in HMDs}
Over the past several decades many studies have found that distances are underestimated in virtual reality (VR) \cite{creem-regehr_chapter_2015, el_jamiy_distance_2019, el_jamiy_survey_2019, renner_perception_2013}. A variety of theories have been proposed to explain this phenomenon, but a comprehensive explanation has not been found \cite{creem2023perceiving}. Proposed theories on the underlying cause of systematic distance underestimation in HMDs are typically related to effects of HMD hardware, such as field of view (FOV) \cite{kline_distance_1996,buck_comparison_2018}, weight \cite{willemsen_effects_2009,buck_comparison_2018,masnadi_field_2021,masnadi_effects_2022}, display resolution \cite{thompson_does_2004,kelly2022distance,ryu_influence_2005}, and higher-order perceptual effects like scene complexity, realism, presence, and embodiment \cite{el_jamiy_distance_2020,bodenheimer_perceiving_2023,vaziri_impact_2017, gonzalez2019individual}. In a meta analysis, Kelly~\shortcite{kelly2022distance} found that distance is more accurately estimated in newer HMDs. Kelly explored the impact of technical characteristics like FOV and resolution on the magnitude of underestimation across studies, but there are other improvements in newer hardware that are more difficult to quantify such as more accurate head tracking, reduced distortion from improved optical designs, and better fitment systems. These advancements all improve the underlying stereo projection and relatively little work has investigated the effect low-level errors in stereo geometry may have on distance perception in HMDs. Our work is the first to explicitly measure how errors in stereo geometry affect behavior (reaching) which may offer new insights on topics like distance underestimation.

\paragraph{Reaching Experiments in Personal Space}
There are fewer studies on distance estimation in VR for personal space (<2 meters) compared to action space (2-30 meters). Results are mixed showing overestimation \cite{rolland_towards_1995}, underestimation \cite{napieralski_near-field_2011}, and accurate estimation \cite{naceri_depth_2011}. Similar to studies in action space, distance is typically measured using open-loop motor tasks such as blind reaching/pointing or verbal reports. Napieralski et. al \shortcite{napieralski_near-field_2011} found the difference between measurement protocal (blind reaching vs. verbal response) is much greater than the difference between VR and the real world (underestimation in both cases). Closed-loop visual feedback has been shown to improve distance judgement accuracy \cite{kelly_recalibration_2014, mohler_influence_2006, altenhoff_effects_2012}, with visual motor re-calibration leading to more accurate reaching \cite{ebrahimi_effects_2014, ebrahimi_carryover_2015}. One challenge in characterizing reach distance in near space is co-registration of the viewer's egocentric coordinate frame and the coordinate frame of the system used to measure reach distance. Our stereo geometry framework in Section~\ref{sec:model} explores how differences in eye relief in particular may introduce discrepancies between how far a user perceives their reach to be and physically measured reach distance by an external observer.

\paragraph{Perceptual Effects of Inaccurate Geometry in HMDs}
Geometric errors have been shown to affect distance perception experimentally, but most studies evaluate engineering-related errors rather than geometric errors explicitly. For example, geometric calibration alleviates distance underestimation in AR \cite{kellner_geometric_2012}. Magnification also causes significant changes in distance judgments \cite{kellner_geometric_2012, kuhl_hmd_2009, zhang_minification_2012, li_effects_2015}. Some studies have examined the impact of mismatched HMD lens interaxial distance (IAD) to the user's IPD and have found that distance perception is minimally impacted by the IAD-IPD mismatch in action space \cite{chakraborty_inter-pupillary_2024,willemsen_effects_2008, hibbard_implications_2020}, likely a result of stereo cues being less accurate at farther distances \cite{mccann_contributions_2018}. \citeN{guan2023perceptual} detection thresholds to eye relief, IPD, and translational viewing errors, but they do not quantify the impact of these thresholds on user behavior or perception. Many computational models have been presented that examine the effects of viewing errors (i.e., displacement of the user CoP from the HMD CoP) \cite{woods93,held2008misperceptions,rolland04}, but fewer have considered the joint effects of simultaneous viewing and rendering errors \cite{holloway1997}. In this work, we complement prior geometric modeling by also pairing geometric predictions with empirical user studies.

%, but not incorrectly calibrating lens distortion\cite{kuhl_hmd_2009}.
% In a theoretical analysis, \cite{hibbard_implications_2020} showed a mismatched interpupillary distance (IPD) would cause inaccurate depth judgement, particularly that smaller (larger) IPD leads to overestimation (underestimation). 

% Here we propose a first-principles based framework that is applicable to all egocentric display architectures relying on stereoscopic perspective projection.

% \paragraph{Models of Geometric Error in Stereoscopic Displays}
% HMD render and viewing geometry assumes that the render cameras, center of projection (CoP) of the displays, and the user's eyes are all co-located in order to render, present, and perceive accurate perspective 3D geometry \cite{shirley2009fundamentals}. Rendering errors can occur when the headset position used for rendering is incorrect and the render camera positions are displaced from the display CoPs (e.g., due to movement of the headset during the time it takes to render and present images, i.e., latency). Viewing errors can occur when the user's eyes are not located at the nominally designed position within the headset (e.g., due to differences in how headsets fit) \cite{woods93,held2008,rolland04}. Fewer studies have separated out viewing and rendering errors \cite{holloway1997}.

% \paragraph{Perceptual Consequences of Render and Viewing Errors}
 % Viewing errors have not been thoroughly studied in the context of visual discomfort, but some studies have examined their potential perceptual implications with psychophysics in headset contexts \cite{kelly13,guan2023perceptual,krajancich2020}. Conversely, rendering errors resulting from latency and lower-accuracy head tracking in older technologies could result in very large errors \cite{holloway1997}, and these types of errors have been more carefully studied in the context of visual discomfort \cite{allison1999}. These dynamic, interactive effects illustrate that correctly making a single estimate of distance doesn't mean that the scene is free of perceptual distortions. They are beyond the scope of this work, but are important to note.